\begingroup
\setcounter{figure}{0}
\renewcommand{\thefigure}{结.\arabic{figure}}
\chapter*{全书结语：从一次操作到可以延续的业务历史}
\phantomsection\label{book:closing}
\addcontentsline{toc}{chapter}{全书结语：从一次操作到可以延续的业务历史}

\begin{center}
  \includegraphics[width=\linewidth,height=.35\textheight,keepaspectratio]{assets/chapter-art/book-closing.png}
  \par\vspace{4pt}
  {\scriptsize\color{Muted}取意于卡斯帕·大卫·弗里德里希《生命的阶段》\par}
  \BookEpigraph{closing}
\end{center}

Alice 再次打开购买页面时，系统可能已经换过工作者、切过副本，也把一部分数据移到了新的节点。她关心的事情依然具体：那次确认有没有付款，买到的是哪份订单，货物走到了哪里，退款有没有回来。回答这些问题的记录分散在购买权威、支付平台和履约系统中，各自经历过不同的等待与恢复。页面把它们接回同一份购买，才有了一段能够连续理解的经历。

最初的故障只隔着一次没有返回的响应。账户已经从 10000 分变成 7000 分，调用方却看不到提交，第二次尝试又扣去 3000 分。沿着这个缺口，设计逐步增加了身份、仲裁、共同提交与原结果。走出本地事务以后，付款可以早于本地确认，消息可以早于发布标记，履约又可以早于消费确认。系统随后还要更换承载位置、缩减历史、恢复备份，并在资源有限的部署中处理新业务与旧责任。

现在可以把这些机制放在一起，重新看一项业务怎样成立。用户的一次确认跨过多次执行，局部决定通过保存的关系组成完整服务，后继执行者从已有证据接续工作。数据结构、组件协议与运行安排各自影响这段历史。设计过程也由此有了次序：先说明业务允许什么，再决定哪些位置保存和仲裁，随后安排历史延续与实际运行。

\section*{一项业务怎样跨越多次执行}
\pdfbookmark[1]{一项业务怎样跨越多次执行}{closing-identity}

同一份输入可以来自两次合法购买，也可以来自一次购买的两次尝试。用户确认的次数决定业务意图，网络发送的次数描述通信活动。调用方在确认时保存完整身份与输入，后续尝试才有共同的起点。本书使用的账户、动作、会话与键，各自确定一层范围。账户联系授权对象，动作解释业务规则，会话带入原有准入条件，键区分该范围内的确认。

这份身份进入服务以后，还要对应实际发生的业务。已有提交绑定指向哪份购买、付款责任与原结果，由权威中的记录决定。一次尝试在提交前回滚，后来可能由另一尝试先形成绑定。提交已经成立，后续同参恢复沿原关联继续。原绑定金额为 3000 分，后续同身份改成 3500 分时，服务返回冲突，保存的 P/C/T 和原金额保持原来的关联。

与此同时，业务约束检查这项购买本身是否合法。两个不同身份都想使用同一账户余额，它们应分别进入资金判断。相同身份的两个请求同时到达，它们则要在相同受理位置汇合。账户约束、身份约束、实际订单与流水的关联共同决定结果。唯一行数只是其中一项关系，金额、对象、授权和返回内容也各有对应事实。

因此，设计一项新服务时，可以先把允许行为写清楚：新确认在什么条件下成功或拒绝，原确认怎样恢复，更换输入怎样处理，查询报告哪个对象和哪个时刻。随后再安排字段、索引和事务。这样才能判断某个技术选择是否保存了业务真正需要的区别。例如，把全部请求按报文摘要合并，会把两次相同金额的合法购买混在一起；每次重试重新生成身份，又会让同一次确认取得新的首次执行机会。

原结果与当前状态进一步说明了观察的含义。在独立的本地购买历史中，Alice 从 10000 分支付 3000 分，当次结果保存 7000 分。随后另一项合法支出使余额变成 6000 分。恢复第一次购买取得当时的 7000 分，当前账户查询报告 6000 分。两个数分别回答不同问题，后续支出沿自己的身份与流水解释。到了外部付款系统，同样的区分落在原收款、原订单，以及当前履约与退款进度上。

用户取消一次等待，也要放在这些关系中理解。关闭页面可以结束一条连接，持久受理的付款责任可能已经存在。申请撤销则是一项业务动作，权威根据原状态判断能否终止尚未发送的工作，或接入原付款核查与退款流程。页面分别显示连接结束、撤销申请和原购买进度，后台再把撤销的处理结果接回原操作。

第十三章讨论的状态表达，为设计提供了另一种选择。标签按集合并入，同一成员多次传播可以汇合，资金增量重复到达则可能再次改变金额。即使一种操作连续重复后状态保持，另一合法操作插入以后也可能产生不同后果。判断幂等、交换、单调和合并收敛时，所比较的对象、允许的交错与业务不变量一起进入分析。标签可以合并重复成员，退款则核对原收款和剩余额度，再为这次退款保存独立占用。

\section*{多个局部决定怎样组成完整服务}
\pdfbookmark[1]{多个局部决定怎样组成完整服务}{closing-cooperation}

第一部分把身份、扣款、订单和原结果放进同一事务。这里有一个能够决定整组修改的位置，提交前退出就共同撤销，提交完成则共同恢复。选择这种结构时，关联写入越容易落在同一事务域，推理与恢复路径越集中。分区设计也由此从业务关系开始：哪些事实经常共同改变，哪些共享资源会形成竞争，怎样安排它们的归属。

支付平台加入以后，购买受理与实际收款有了各自的提交点。平台已经完成 X1、本地尚未回填，和请求尚未到达平台，都可能留下相同的本地处理中状态。恢复者从原 C 和平台范围进入查询，依据可信响应、回调或对账确定发生过什么。平台对原键的仲裁能力、结果可查范围与保护期限，直接改变允许的发送动作。

原身份把多次尝试指向相同工作，接收端的仲裁决定这些尝试怎样形成实际效果。购买库保存一条 C，解释本地已经承担哪项责任，平台用原范围与原键限制自己的收款。一次租约接管让新工作者安排下一轮查询，已经离开旧工作者的请求则沿平台操作关联恢复。平台结果进入来源与内容核验，资源修改则同时检查代次与当前状态，再提交相应变化。

跨域接续还需要一个能够在进程退出后重新发现的保存位置。受理提交先保存原输入、外部关联和 T，付款结果提交再保存订单、原结果与 EV。消息发布器从 EV 继续传送，消费者把接收、业务守卫、履约与发运意图共同提交。各段的完成证据再返回原责任。前一个进程已经结束，后一个执行者依然能从确定的位置找到工作。

图~\ref{fig:closing-history} 选取最终系统中已满足付款资格的一份购买，沿一条成功付款、随后进入履约与退款的历史，分别列出事实、接续责任和查询含义。

\begin{figure}[htbp]
\centering
\begin{tikzpicture}[
  x=1mm,y=1mm,
  cell/.style={draw=Ink!35,fill=Paper,rounded corners=1.5pt,
    text width=43mm,minimum height=15mm,align=center,inner sep=2mm,
    font=\footnotesize\sffamily},
  fact/.style={cell,draw=Accent!65,fill=Accent!6},
  head/.style={font=\small\sffamily\bfseries,text=Ink,align=center},
  closingflow/.style={-{Stealth[length=1.6mm]},draw=Ink!60,line width=.6pt}
]
\node[head] at (0,5) {已经保存的事实};
\node[head] at (53,5) {接续责任};
\node[head] at (106,5) {查询所报告的含义};
\node[fact] (a1) at (0,-11) {原购买已受理\\P1/C1/T1 已保存};
\node[cell] (b1) at (53,-11) {已满足付款资格\\按原平台关联推进};
\node[cell] (c1) at (106,-11) {已受理，付款尚待确认};
\node[fact] (a2) at (0,-33) {平台 X1 已收款 3000 分\\本地尚缺确定证据};
\node[cell] (b2) at (53,-33) {沿原 C1 查询、回调或对账\\继续核对 X1};
\node[cell] (c2) at (106,-33) {本地报告付款未知\\平台收款事实已经存在};
\node[fact] (a3) at (0,-55) {原结果、O1 与 EV1\\已在权威共同保存};
\node[cell] (b3) at (53,-55) {付款责任结束\\原 EV1 进入传播};
\node[cell] (c3) at (106,-55) {原购买成功\\履约进度另查};
\node[fact] (a4) at (0,-77) {消费者已提交 F1\\并保存交接与发运意图};
\node[cell] (b4) at (53,-77) {交接证据返回上游\\发运沿原身份接续};
\node[cell] (c4) at (106,-77) {履约已接受\\实际发运按其证据报告};
\node[fact] (a5) at (0,-99) {新 RF 关联原 X1\\保存退款及相应占额};
\node[cell] (b5) at (53,-99) {退款从原 RF 继续\\取得确定证据后更新};
\node[cell] (c5) at (106,-99) {原购买结果保持\\当前退款进度继续变化};
\foreach \r in {1,2,3,4,5}{
  \draw[closingflow] (a\r.east)--(b\r.west);
  \draw[closingflow] (b\r.east)--(c\r.west);
}
\foreach \old/\new in {1/2,2/3,3/4,4/5}{
  \draw[closingflow] (a\old.south)--(a\new.north);
}
\end{tikzpicture}

\caption{同一购买的事实、恢复责任与用户观察分别推进，后续阶段沿原关联接续。}
\label{fig:closing-history}
\end{figure}

代理确认描述消息域已经接受什么，消费事务描述履约系统已经保存什么，实际发运又属于外部平台的事实。消息保留期有限时，购买权威根据持久消费交接证据结束原事件的补发责任。消费者已经形成 F，原 dispatch 还可以处于发运未知，发运工作者沿自己的持久记录继续取得外部结果，再更新这份履约的进度。

有些业务关系需要先取得共同决定。组合购买要求两个库存分区联合预留时，参与者保存准备占用，协调者保存唯一决定，再把决定传播到各参与者。已经准备的资源根据可信决定提交或释放。参考流程在必需参与者完成以后开启原 C/T 的付款资格。此时共同决定的范围覆盖联合库存，后续支付继续遵守外部平台的提交与恢复规则。

允许阶段性可见的业务可以采取另一种组织：先完成一个资源预留，再推进另一资源，失败以后沿已经完成的阶段发起补偿。需要联合预留，就承担共同决定、等待占用与恢复的成本，允许部分阶段先成立，就在相应状态中保存已完成的预留、后续动作和补偿身份，让查询与恢复都能识别业务当前处于哪个阶段。

退款说明后续修正怎样进入历史。原收款 3000 分，成功退款 1000 分，另有 1500 分结果未知，后续最多再占用 500 分。未知退款转给人工核查以后，原 RF 与占额继续保存。取得成功或足够最终拒绝的证据，权威再转换或释放占用。原购买成功、当前净资金和退款进度分别由关联事实解释，修正动作取得自己的授权、身份与结果。

\section*{系统怎样把过去交给后来者}
\pdfbookmark[1]{系统怎样把过去交给后来者}{closing-continuity}

前面的每一段恢复，都使用了某种已经保存的区别：这个身份已经受理，那份收款属于原 C，这次履约已经形成，某项退款还占用额度。复制、快照、回收与迁移改变这些区别的承载方式。设计要追踪后继者取得了哪些判断依据，以及这些依据允许它采取什么动作。

复制首先保护已经确认的决定。状态机按提交顺序应用命令，把业务关系、会话、结果、编号进度和未决责任保存在状态中。接班副本从确定的前缀继续，后来的重复命令找到原身份。数据库事务也可以通过内部复制保存提交结局。判断一个成功响应经历故障后怎样延续，要沿实际确认路径核对复制、持久化与应用位置。

事务与复制由此进入同一条业务路径，但各自说明不同关系。事务决定哪些变化共同成立，复制说明相应决定怎样跨节点故障保存。跨多个复制组的资源修改，还需要共同决定协议。一个系统可以同时具有组内共识、跨组事务和外部工作流，它们按各自范围衔接，恢复者分别从日志前缀、事务决定和外部操作证据继续。

快照把一段历史压缩为某个应用位置的完整状态。只复制账户和订单，可能遗漏解释旧请求的身份绑定，或让尚未结束的付款失去继续入口。完整图像把恢复所需的关系与边界放在一起。日志截断发生以后，后继者由图像取得这些区别，再应用后缀命令。图像格式、发布位置和校验信息也因此成为恢复协议的一部分。

回收进一步提出一个问题：逐项结果删除以后，迟到请求得到什么处置？关闭会话、保留必要拒绝依据、限制原提交入口，是书中采用的安排。保存方式可以从详细结果变成较小的关闭证据，旧入口据此作出确定判断。已经受理的 C/T 则通过后台恢复与授权查询继续。客户端重新提交的资格和服务已经承担的责任，按各自结束条件处理。

这项选择把存储成本与业务能力联系起来。保留详细结果，能够在相应范围内直接恢复原响应；缩减结果，查询与拒绝规则随之改变。外部平台的原键窗口和在途请求还会限制责任如何终结。单独减少本地保留时间，可能只是让恢复者更早失去证据。确定回收策略时，要把入口期限、外部效果与服务端持续责任一同安排。

在线迁移改变完整业务分区的承载位置。源组先形成确定的冻结边界，目标导入图像，目录再发布新归属与资格。参考首版在 prepared 库存责任或协调者事务未闭合时等待相应决定完成，然后发起冻结。已经受理的外部付款作为持久责任随图像交接，目标从原绑定取得相同的 C 与平台键，继续查询先前发送的请求。

底层 Region 调度、Redis 槽迁移和完整业务分区迁移，有不同的对象与完成条件。一个区间或槽搬走以后，跨范围的业务关系由上层事务、路由与恢复安排维护。上层把原绑定、订单、事件与未决任务的恢复关系列出来，再核对这些对象在底层迁移中的位置和可读时刻。

旧备份恢复更接近开篇的地域故事。本地历史退回，外部收款与履约可能继续保留较新的事实。新的执行代次可以保护新编号与执行资格，旧购买却需要找回自己的原绑定。独立保存的关联和平台账本让恢复者重建原 P/C/T、订单与事件，再按完整恢复的范围开放入口。歧义范围通过持久关闭限制新提交，新的业务获得清楚的准入边界。

升级同样要把过去交给新代码。动作版本决定旧输入怎样解释，原结果版本决定历史响应怎样读取，未完成工作带着原关联进入后继实现。改变业务规则以后，新确认按新规则判断，旧身份的恢复按原绑定继续。新代码先取得解释旧记录所需的格式和版本，再接管原 C/T，依照保存的输入、平台范围与结果推进后续步骤。

\section*{基础设施怎样改变实现方式与成本}
\pdfbookmark[1]{基础设施怎样改变实现方式与成本}{closing-infrastructure}

业务关系、决定范围与延续条件确定以后，就可以为它们选择具体的承载方式。消息系统承担传播，数据库或状态机承担权威修改，控制面管理归属与资格，缓存形成便于读取的观察，存储与计算组件保存和执行大规模工作。

RabbitMQ 围绕路由、队列与确认组织工作交付，Kafka 围绕分区日志、保留和消费位置组织读取。它们分别适合不同的消费与重读安排。购买事件仍以原 EV 关联业务，消费者在自己的效果位置检查接收与业务守卫。选型时，事件传播量、重读需求、分区顺序、消费并发和保留时间共同决定消息域的结构，消费事务与交接证据再把它接回完整服务。

缓存把读取从业务权威移到派生状态。用于页面观察时，可以依据用途安排刷新、失效与落后提示；用于判定新购买能否受理时，读取必须满足该项判断需要的事实条件。Redis 的命令执行、RDB、AOF、复制与切换规则，具体决定状态怎样进入保存与恢复路径。把某项决定交给它承担，就沿这些路径核对相应保证。把它用于可重建缓存，则把来源、版本和重建预算保存清楚。

Raft 和链式复制安排副本如何共同保存状态，ZooKeeper、etcd 或 Consul 提供管理成员、会话、版本和资格的接口。取得资格以后，实际资源在写入处核验它，再提交相应修改。路由缓存降低对控制面的访问。以 etcd 目录为例，通知中断后，入口根据已保存的完整修订位置接续事件，或在历史已压缩时重新取得完整基线，再更新本地路由。

存储层把另一项决定显露出来：候选字节何时成为有效状态。块和文件服务保存内容、版本与元数据，对象接口提供自己的条件操作。完整图像、输出文件集合和湖仓快照则经过各自的发布步骤。读者固定已发布的版本或集合，恢复者据发布事实处理候选与旧文件。对象已经存在、当前集合已经更新、应用已经取得持久确认，分别位于这条路径上的不同位置。

作业和运行时把执行次数带到同一问题上。MapReduce、Spark 等可以重新执行或推测执行任务，输出提交决定哪个结果进入有效集合。Ray 的任务、对象与 Actor 有各自的重建和重试条件，函数运行环境也可能替换或复用。原业务身份随工作进入后继执行环境，外部效果通过保存的关联核对。持久工作流保存步骤位置以后，付款步骤通过原 C 与原键恢复平台事实。

前沿方案进一步移动等待和保存的位置。推测执行让内部计算与后台持久化重叠，外部效果依据相应屏障取得继续条件；共识快路径改变决定的通信安排；地域读协议把安全读取放到选定响应者；对象存储引擎与表格式把有效状态发布到新的承载层。比较这些方案，可以追踪哪一段等待缩短、增加了哪些通知或元数据，以及故障后谁重新取得确定位置。

选择状态表达与部署方式时，业务允许合并的状态可以减少协调，共享资金、库存或额度的首次承诺则落在能够仲裁资源的位置。把互不重叠的额度预分配到地域，可以让当地在获配范围内继续受理，委托使用与失联结算成为新增责任。把相关数据放近可以减少跨域决定，热门账户与 SKU 的请求又可能集中竞争相同资源，容量与准入据这些共享点分别安排。

图~\ref{fig:closing-design} 将这种设计次序连起来。实现选择完成以后，验证与运行会提供新的信息。负载分布变化可能要求重新分区，外部平台能力变化可能要求调整恢复路径，新的业务动作则先回到允许行为与资源约束，再沿后续关系检查。

\begin{figure}[htbp]
\centering
\begin{tikzpicture}[
  x=1mm,y=1mm,
  box/.style={draw=Ink!45,fill=Paper,rounded corners=2pt,
    text width=42mm,minimum height=22mm,align=center,inner sep=2mm,
    font=\small\sffamily},
  chosen/.style={box,draw=Accent!70,fill=Accent!7},
  next/.style={-{Stealth[length=1.8mm]},draw=Ink,line width=.7pt},
  back/.style={-{Stealth[length=1.8mm]},draw=Accent,line width=.7pt}
]
\node[chosen] (meaning) at (0,0) {业务允许行为\\身份、授权、资源约束\\原结果与当前观察};
\node[box] (decision) at (53,0) {事实与决定的归属\\共同提交范围\\跨域责任与证据};
\node[box] (history) at (106,0) {后继者的恢复依据\\复制、回收、迁移\\版本与旧历史恢复};
\node[box] (implementation) at (106,-35) {承载与实现选择\\存储、消息、计算\\资格与发布协议};
\node[box] (validation) at (53,-35) {论证与历史检查\\不变量、交错与故障\\核对实际保存事实};
\node[chosen] (operation) at (0,-35) {负载与长期运行\\准入、调度、排空\\观测与恢复预算};
\draw[next] (meaning.east)--(decision.west);
\draw[next] (decision.east)--(history.west);
\draw[next] (history.south)--(implementation.north);
\draw[next] (implementation.west)--(validation.east);
\draw[next] (validation.west)--(operation.east);
\draw[back] (operation.north)--node[left,font=\scriptsize\sffamily,text=Accent]
  {条件变化} (meaning.south);
\end{tikzpicture}

\caption{业务要求先确定需要保存和仲裁的关系，实现、验证与运行再共同检验这些选择。}
\label{fig:closing-design}
\end{figure}

\section*{正确的设计怎样持续运行}
\pdfbookmark[1]{正确的设计怎样持续运行}{closing-operation}

一份持久责任即使保存完整，也要得到实际执行机会。数据库恢复提交能力，旧事务结束占用，工作者取得调度，网络给出通信机会，平台最终提供能够核验的结果，这些条件分别使某一段工作得以继续。业务提交以后，客户端还要经过结果读取、响应传输与本地保存，才能在下一次启动时直接找到已经确认的购买。

新购买、同身份重试、未知付款查询、事件发布和消费恢复，使用的是不同工作单位。一次重试可以汇合到原业务效果，连接、带宽和平台调用额度已经消耗。入口合并观察减少重复读取，后台按到期时间安排查询，退避分散突发，有界并发限制实际在途压力。平台额度按尝试计量，责任排空按确定完成的业务计量，查询取得确定证据的比例把两者接到一起。

沿用第二十一章的排空例子，各阶段能力在计算窗口内稳定，每份付款终结产生一份事件。平台暂停十分钟，每秒受理 200 份，形成 12 万份 T。恢复后每秒终结 500 份、新流继续每秒 200 份，付款积压约需 400 秒排空。这段时间每秒产生约 500 份 EV，消费能力若只有每秒 350 份，又会累积约 6 万份消费责任。付款恢复以后，消费再以每秒 150 份的净速度处理，整条链约在第 800 秒完成排空。

因此，放量依据整条链的余量安排。新受理会增加未来的付款、事件与履约责任，确定结果回填则可能立即结束一份未知。原结果查询、事务决定恢复、事件交接和旧责任查询各自取得预算，暂时空闲的份额再按规则借用。某个分区的热点受到批次、并发与配额约束，其他分区的少量未决工作也能获得处理机会。

公共身份资格检查通过以后，已有绑定沿原记录恢复，新身份再参与新增责任额度的仲裁。会话关闭以后，后台责任通过自己的入口继续。忙、尚未受理、结果未知和确定业务拒绝分别对应具体的受理状态，调用方据此保留原身份查询、等待下一次机会，或重新确认经过修改的新购买。

验证从这些允许行为和迁移关系开始。给定明确初态，列出请求、提交、丢失、接管和恢复的先后，分别记录各参与者能观察什么，再与独立规格比较。证明检查允许迁移怎样保持不变量，有界模型展开指定范围的交错，故障实验观察实际实现留下什么，恢复演练把备份、路由、凭据、外部证据和容量接进整条链。

这些方法还可以指导问题定位。出现两次实际收款，沿原身份、平台范围和首次仲裁查找；同一笔收款长期无法结束，检查证据取得、关联核验与结果回填；迁移后旧请求成为新购买，追踪图像、准入边界和原绑定。观测保留同一业务的关系，反例就能落到具体状态与步骤，随后修改相应设计并重新检查受影响的历史。

系统演进会不断改变这些检查的输入。加入充值以后，余额关系同时计入收入与支出，原先依赖余额只降的拒绝判断重新进入并发历史。允许跨地域受理以后，归属、额度与仲裁按新的入口安排。更换平台以后，原键窗口与查询能力进入新的故障历史，已经受理的旧 C 则保留原平台关联，由兼容的恢复路径继续核查。

\section*{回到最初的交接现场}
\pdfbookmark[1]{回到最初的交接现场}{closing-return}

实验室里的 Alice 面对一份不能重新取得的样本，操作单上“加入试剂”一栏空白。Bob 已经加过却未登记，和 Bob 尚未动手，可以留下相同的记录。我们现在能够更准确地描述她的处境：眼前记录缺少区分这两条历史的信息，重复操作的后果却由样本实际经历决定。她暂停动作、保留现场，等待 Bob 的确认或可靠检测，把下一步交给能够取得的证据。

计算机系统把一部分缝隙收进原子提交。身份、余额、订单与原结果能够共同保存时，接班程序从同一结局恢复。独立支付、实际发运和试剂加入则有各自的效果位置。设计先确定这些位置能提供什么仲裁和查询，再根据原操作可取得的证据，安排重发、核查或持续等待的条件。

Bob 的迟来消息涉及执行安排和已发生事实。他说“已经加过”，带来的是需要核对对象与步骤的证据，他说“现在去加”，则涉及交接以后的执行资格。购买系统中的可信旧平台结果由服务核对原 C、平台范围与完整输入，迟到工作者的资源写入则根据当前代次和状态取得处理结果。

游戏中的甲地域则恢复了另一份过去。旧角色、未消耗材料和待迁移任务在备份里彼此一致，乙地的活动与 Alice 持有的装备却已经向前。恢复者从原迁移身份找到完成记录，核对后来转移的归属，再决定哪些对象能够恢复、哪些入口可以开放。旧备份描述当时保存过什么，当前处置权由能够延续这段历史的记录与规则确定。

若关联记录已经回收，恢复范围和剩余证据就进入判断。关闭旧入口可以阻止歧义任务重新取得资格，可靠外部历史可以帮助重建原关系，人工核查则保留未决对象与处置责任。恢复者根据业务允许的等待、当前可以取得的证据和故障前实际保存的关联，确定哪些旧工作能够继续核查，以及哪些范围可以重新开放。

最后回到 Alice 的购买。原身份找到当时形成的订单与收款，查询再显示当前履约和退款。某个阶段尚未结束，持久记录说明最近取得了什么证据、谁继续办理、下一步安排在何处。换一个执行者、换一组副本，接管者从这些关系进入已经发生的历史，把尚未结束的工作继续下去。

\FloatBarrier
\clearpage
\endgroup
